\chapter[Asterisk instalazioa]{Asterisk instalazioa}
\label{chp:asterisk-install}

\localtableofcontents

\section{Sarrera}

Asterisk softwareak gure zoiper bezeroei zerbitzua emango die. Instalazio hau ere Docker irudiak erabiliz egingo da.

\section{Instalazioa}

Hurrengo pausoak jarraituz Asterisk aplikazioa instalatuko da Docker edukiontzi batean.
\begin{enumerate}
    \item Karpeta prestatu instalazioa egiteko: \mintinline{text}{~/Documents/VoIP/Asterisk/}
    \begin{minted}{console}
    $ mkdir -p ~/Documents/VoIP/Asterisk
    \end{minted}
    \item Asterisk-en kodea deskargatu aurretik sortutako karpetan.
    \begin{minted}{console}
    $ cd ~/Documents/VoIP/Asterisk
    $ wget http://downloads.asterisk.org/pub/telephony/asterisk/asterisk-20-current.tar.gz
    \end{minted}
    \item Fitxategia deskonprimitu.
    \begin{minted}{console}
    $ tar -zxvf asterisk-20-current.tar.gz
    \end{minted}
    \item \chapterref{chp:zoiper-install}-en sortutako Docker irudia exekutatu.
    \begin{minted}{console}
    $ docker run -it -v /home/$USER/Documents/VoIP/Asterisk:/Asterisk ubuntu:zoiper
    \end{minted}
    \item Hurrengo pausoak jarraitu Asterisk konpilatzeko iturburutik.
    \begin{minted}{console}
    # cd /Asterisk/asterisk-20.10.0
    # cd contrib/scripts
    # ./install_prereq install
    # cd ../..
    # ./configure
    # make
    # make install
    # make samples
    \end{minted}
    \item \textbf{Docker terminaletik kanpo}, gorde aldaketak irudi berri batean honako komandoa exekutatuz:
    \begin{important}
        ID egokia emango du baldin eta exekutatu zen azkeneko docker edukiontzia asterisk-ena bazen.
    \end{important}
    \begin{minted}{console}
    $ newContainerID=$(docker ps -a | awk 'NR==2{printf $1}') && docker commit $newContainerID ubuntu:asterisk
    \end{minted}
\end{enumerate}

\section{Exekuzioa}

Hurrengo pausoek Asterisk exekutatzeko pausoak adierazten ditu.
\begin{enumerate}
    \item Docker edukiontzia exekutatu hurrengo agindua erabiliz:
    \begin{minted}{console}
    $ docker run -it -v /home/$USER/Documents/VoIP/Asterisk:/Asterisk ubuntu:asterisk
    \end{minted}
    \item Asterisk exekutatu edukiontzian:
    \begin{minted}{console}
    # asterisk -cvvvvv
    \end{minted}
\end{enumerate}

Aldaketak egindakoan asterisk-en konfigurazio fitxategiei, \textbf{aldaketak gordetzeko} hurrengo komandoa exekutatu kontsola berri batean:
\begin{important}
    ID egokia emango du baldin eta exekutatu zen azkeneko docker edukiontzia asterisk-ena bazen.
\end{important}
\begin{minted}{console}
$ newContainerID=$(docker ps -a | awk 'NR==2{printf $1}') && docker commit $newContainerID ubuntu:asterisk
\end{minted}